\begin{tikzpicture}[
  x=1mm,y=1mm,
  >=Stealth,
  font=\figurefont\fontsize{8}{10}\selectfont
]

\tikzset{
  vnode/.style={circle,draw=FigureBlue,fill=FigureBlueFill,minimum size=3.8mm,inner sep=0pt,line width=.5pt,text=FigureInk,font=\figurefont\fontsize{6.8}{8}\selectfont},
  vstart/.style={circle,draw=FigureYellow,fill=FigureYellowFill,minimum size=4.0mm,inner sep=0pt,line width=.55pt,text=FigureInk,font=\figurefont\fontsize{6.8}{8}\selectfont},
  vtarget/.style={circle,double,draw=FigureRed,fill=FigureRedFill,minimum size=4.0mm,inner sep=0pt,line width=.5pt,text=FigureInk,font=\figurefont\fontsize{6.8}{8}\selectfont},
  edge/.style={draw=FigureStroke,line width=.5pt},
  note/.style={font=\FigureNoteFont\fontsize{7.2}{8.8}\selectfont,text=FigureInk},
  muted/.style={font=\FigureNoteFont\fontsize{6.8}{8.2}\selectfont,text=FigureMuted}
}

% Row 1, Col 1: (a) Complete graph K_5
\begin{scope}[xshift=0mm,yshift=26mm]
  \node[vstart] (k1) at (3,12) {$i$};
  \node[vnode] (k2) at (12,15) {};
  \node[vtarget] (k3) at (21,12) {$j$};
  \node[vnode] (k4) at (18,3) {};
  \node[vnode] (k5) at (6,3) {};
  \foreach \u/\v in {k1/k2,k1/k3,k1/k4,k1/k5,k2/k3,k2/k4,k2/k5,k3/k4,k3/k5,k4/k5}{
    \draw[edge] (\u)--(\v);
  }
  \node[note] at (12,-1.5) {(a) 完全图 $K_n$};
  \node[muted] at (12,-5.5) {$h_{i,j}=n-1=\Theta(n)$};
\end{scope}

% Row 1, Col 2: (b) Star graph K_{1,4}
\begin{scope}[xshift=35mm,yshift=26mm]
  \node[vstart] (s0) at (12,8) {$c$};
  \node[vnode] (s1) at (3,14) {$u$};
  \node[vnode] (s2) at (3,2) {};
  \node[vtarget] (s3) at (21,14) {$v$};
  \node[vnode] (s4) at (21,2) {};
  \draw[edge] (s0)--(s1);
  \draw[edge] (s0)--(s2);
  \draw[edge,draw=FigureRed,line width=.7pt] (s0)--(s3);
  \draw[edge] (s0)--(s4);
  \node[note] at (12,-1.5) {(b) 星图 $K_{1,n-1}$};
  \node[muted] at (12,-5.5) {$h_{u,c}=1$，$h_{c,v}=2n-3$};
\end{scope}

% Row 1, Col 3: (c) Path graph P_n
\begin{scope}[xshift=70mm,yshift=26mm]
  \node[vstart] (p1) at (0,8) {$1$};
  \node[vnode] (p2) at (8,8) {$2$};
  \node[vnode] (p3) at (16,8) {};
  \node[vtarget] (p4) at (24,8) {$n$};
  \draw[edge] (p1)--(p2)--(p3)--(p4);
  \node[note] at (12,-1.5) {(c) 路径图 $P_n$};
  \node[muted] at (12,-5.5) {$h_{1,n}=(n-1)^2=\Theta(n^2)$};
\end{scope}

% Row 2, Col 1: (d) Cycle graph C_n
\begin{scope}[xshift=0mm,yshift=0mm]
  \node[vstart] (c1) at (3,9) {$i$};
  \node[vnode] (c2) at (8,15) {};
  \node[vnode] (c3) at (16,15) {};
  \node[vtarget] (c4) at (21,9) {$j$};
  \node[vnode] (c5) at (16,3) {};
  \node[vnode] (c6) at (8,3) {};
  \draw[edge] (c1)--(c2)--(c3)--(c4)--(c5)--(c6)--(c1);
  \node[note] at (12,-1.5) {(d) 环图 $C_n$};
  \node[muted] at (12,-5.5) {$h_{i,j}=k(n-k)=\Theta(n^2)$};
\end{scope}

% Row 2, Col 2: (e) Hypercube Q_3
\begin{scope}[xshift=35mm,yshift=0mm]
  \node[vstart] (q1) at (3,5) {$i$};
  \node[vnode] (q2) at (13,5) {};
  \node[vnode] (q3) at (3,13) {};
  \node[vnode] (q4) at (13,13) {};
  \node[vnode] (q5) at (8,7.5) {};
  \node[vnode] (q6) at (18,7.5) {};
  \node[vnode] (q7) at (8,15.5) {};
  \node[vtarget] (q8) at (18,15.5) {$j$};
  \draw[edge] (q1)--(q2)--(q4)--(q3)--(q1);
  \draw[edge] (q5)--(q6)--(q8)--(q7)--(q5);
  \draw[edge] (q1)--(q5);
  \draw[edge] (q2)--(q6);
  \draw[edge] (q3)--(q7);
  \draw[edge] (q4)--(q8);
  \node[note] at (12,-1.5) {(e) 超立方体 $Q_d$};
  \node[muted] at (12,-5.5) {多路分流 $h_{i,j}=\Theta(n)$};
\end{scope}

% Row 2, Col 3: (f) Lollipop graph
\begin{scope}[xshift=70mm,yshift=0mm]
  \node[vstart] (l1) at (0,13) {};
  \node[vnode] (l2) at (0,5) {};
  \node[vnode] (l3) at (7,13) {};
  \node[vnode] (l4) at (7,5) {};
  \node[vnode] (l5) at (12,9) {$u$};
  \node[vnode] (l6) at (19,9) {};
  \node[vtarget] (l7) at (26,9) {$v$};
  \foreach \u/\v in {l1/l2,l1/l3,l1/l4,l1/l5,l2/l3,l2/l4,l2/l5,l3/l4,l3/l5,l4/l5}{
    \draw[edge] (\u)--(\v);
  }
  \draw[edge,draw=FigureRed,line width=.65pt] (l5)--(l6)--(l7);
  \node[note] at (13,-1.5) {(f) 棒棒糖图};
  \node[muted] at (13,-5.5) {团内困滞 $h_{u,v}=\Theta(n^3)$};
\end{scope}
\end{tikzpicture}
